% GENERATED FILE. Edit canonical lessons, then run npm run generate:books.
% canonical-source-hash: fnv1a64:933334c0239d7aba
% canonical-lessons: AR-C398-talim, AR-C398-tarbiya, AR-C398-rawda, AR-C398-mahw-al-ummiyya, AR-C398-hudur

\chapter{Schooling}
\label{ch:ar-schooling}
\hlchaptermodality{398}

\begin{chapteropening}
\textbf{By the end of this chapter:} \emph{I can talk about education, upbringing, a kindergarten, literacy and attendance.}

Complete the last lesson of chapter 398: I can talk about education, upbringing, a kindergarten, literacy and attendance.
\end{chapteropening}

\section[\texorpdfstring{\ar{تعليم} (taʿlīm)}{تعليم (taʿlīm)}]{\ar{تعليم} (taʿlīm) — education}

\label{lesson:AR-C398-talim}



\begin{quote}
Before the new one: say the Arabic for a syllabus, then the Arabic for training.
\end{quote}

\subsection*{You'll want to know: \ar{تعليم}}
\textbf{\ar{تعليم}} — \emph{taʿlīm} — \textquotedblleft{}education\textquotedblright{}.

\begin{quote}
\ar{التعليم} \ar{مجاني} \ar{في} \ar{هذا} \ar{البلد}.
\end{quote}

\begin{quote}
\emph{at-taʿlīm majjānī fī hādhā al-balad.}
\end{quote}

\begin{quote}
Education is free in this country.
\end{quote}

\begin{cousinweb}
\textbf{\ar{تعليم}} (\emph{taʿlīm}) is built on the native Arabic root \emph{ʿ-l-m}, \textquotedblleft{}to know\textquotedblright{} (Lane). \emph{Taʿlīm} is the Form II verbal noun, making someone know: teaching, education.
\end{cousinweb}

\begin{grammarlens}[title={What you've built}]
One more word for giving reasons and opinions.
\end{grammarlens}

\subsection*{Your turn}
\begin{itemize}
\raggedright
  \item \emph{Say it:} \emph{taʿlīm}
  \item \emph{Say it:} \emph{taʿlīm}, once more
  \item \emph{Say it:} \emph{at-taʿlīm majjānī fī hādhā al-balad.} — Education is free in this country.
  \item \emph{Recall:} say the Arabic for a syllabus, then the Arabic for training, then say \emph{taʿlīm} again
\end{itemize}

\begin{tcolorbox}[breakable,colback=teal!4,colframe=teal!35!black,title={Before you move on}]
What does \ar{تعليم} mean? (\textquotedblleft{}Education\textquotedblright{}.) Say it once more.
\end{tcolorbox}
\section[\texorpdfstring{\ar{تربية} (tarbiya)}{تربية (tarbiya)}]{\ar{تربية} (tarbiya) — upbringing}

\label{lesson:AR-C398-tarbiya}



\begin{quote}
Before the new one: say the Arabic for training, then the Arabic for education.
\end{quote}

\subsection*{You'll want to know: \ar{تربية}}
\textbf{\ar{تربية}} — \emph{tarbiya} — \textquotedblleft{}upbringing\textquotedblright{}.

\begin{quote}
\ar{تهتم} \ar{أمي} \ar{بتربية} \ar{أطفالها}.
\end{quote}

\begin{quote}
\emph{tahtammu ummī bi-tarbiyat aṭfālihā.}
\end{quote}

\begin{quote}
My mother takes care over bringing up her children.
\end{quote}

\begin{cousinweb}
\textbf{\ar{تربية}} (\emph{tarbiya}) is built on the native Arabic root \emph{r-b-w}, \textquotedblleft{}to grow, to increase\textquotedblright{} (Lane). \emph{Tarbiya} is the Form II verbal noun, making a child grow: upbringing.
\end{cousinweb}

\begin{grammarlens}[title={What you've built}]
One more word for giving reasons and opinions.
\end{grammarlens}

\subsection*{Your turn}
\begin{itemize}
\raggedright
  \item \emph{Say it:} \emph{tarbiya}
  \item \emph{Say it:} \emph{tarbiya}, once more
  \item \emph{Say it:} \emph{tahtammu ummī bi-tarbiyat aṭfālihā.} — My mother takes care over bringing up her children.
  \item \emph{Recall:} say the Arabic for training, then the Arabic for education, then say \emph{tarbiya} again
\end{itemize}

\begin{tcolorbox}[breakable,colback=teal!4,colframe=teal!35!black,title={Before you move on}]
What does \ar{تربية} mean? (\textquotedblleft{}Upbringing\textquotedblright{}.) Say it once more.
\end{tcolorbox}
\section[\texorpdfstring{\ar{روضة} (rawḍa)}{روضة (rawḍa)}]{\ar{روضة} (rawḍa) — a kindergarten}

\label{lesson:AR-C398-rawda}



\begin{quote}
Before the new one: say the Arabic for education, then the Arabic for upbringing.
\end{quote}

\subsection*{You'll want to know: \ar{روضة}}
\textbf{\ar{روضة}} — \emph{rawḍa} — \textquotedblleft{}a kindergarten\textquotedblright{}.

\begin{quote}
\ar{يذهب} \ar{ابني} \ar{إلى} \ar{الروضة} \ar{كل} \ar{صباح}.
\end{quote}

\begin{quote}
\emph{yadhhabu ibnī ilā ar-rawḍa kulla ṣabāḥ.}
\end{quote}

\begin{quote}
My son goes to kindergarten every morning.
\end{quote}

\begin{cousinweb}
\textbf{\ar{روضة}} (\emph{rawḍa}) is built on the native Arabic root \emph{r-w-ḍ} (Lane), where it names a green meadow or garden. Modern Arabic made it the kindergarten, a garden of children.
\end{cousinweb}

\begin{grammarlens}[title={What you've built}]
One more word for giving reasons and opinions.
\end{grammarlens}

\subsection*{Your turn}
\begin{itemize}
\raggedright
  \item \emph{Say it:} \emph{rawḍa}
  \item \emph{Say it:} \emph{rawḍa}, once more
  \item \emph{Say it:} \emph{yadhhabu ibnī ilā ar-rawḍa kulla ṣabāḥ.} — My son goes to kindergarten every morning.
  \item \emph{Recall:} say the Arabic for education, then the Arabic for upbringing, then say \emph{rawḍa} again
\end{itemize}

\begin{tcolorbox}[breakable,colback=teal!4,colframe=teal!35!black,title={Before you move on}]
What does \ar{روضة} mean? (\textquotedblleft{}A kindergarten\textquotedblright{}.) Say it once more.
\end{tcolorbox}
\section[\texorpdfstring{\ar{محو} \ar{الأمية} (maḥw al-ummiyya)}{محو الأمية (maḥw al-ummiyya)}]{\ar{محو} \ar{الأمية} (maḥw al-ummiyya) — literacy}

\label{lesson:AR-C398-mahw-al-ummiyya}



\begin{quote}
Before the new one: say the Arabic for upbringing, then the Arabic for a kindergarten.
\end{quote}

\subsection*{You'll want to know: \ar{محو} \ar{الأمية}}
\textbf{\ar{محو} \ar{الأمية}} — \emph{maḥw al-ummiyya} — \textquotedblleft{}literacy\textquotedblright{}.

\begin{quote}
\ar{تعمل} \ar{الجمعية} \ar{على} \ar{محو} \ar{الأمية} \ar{في} \ar{القرى}.
\end{quote}

\begin{quote}
\emph{taʿmalu al-jamʿiyya ʿalā maḥw al-ummiyya fī al-qurā.}
\end{quote}

\begin{quote}
The association works for literacy in the villages.
\end{quote}

Literally the wiping out of illiteracy; \ar{أمي}, \emph{ummī}, is one who cannot read.

\begin{cousinweb}
\textbf{\ar{محو} \ar{الأمية}} (\emph{maḥw al-ummiyya}) joins \emph{maḥw}, \textquotedblleft{}wiping out\textquotedblright{}, built on the native Arabic root \emph{m-ḥ-w}, and \emph{ummiyya}, \textquotedblleft{}illiteracy\textquotedblright{}, from \emph{ummī}, built on \emph{ʾ-m-m} (Lane): one still as his mother bore him, unable to read.
\end{cousinweb}

\begin{grammarlens}[title={What you've built}]
One more word for giving reasons and opinions.
\end{grammarlens}

\subsection*{Your turn}
\begin{itemize}
\raggedright
  \item \emph{Say it:} \emph{maḥw al-ummiyya}
  \item \emph{Say it:} \emph{maḥw al-ummiyya}, once more
  \item \emph{Say it:} \emph{taʿmalu al-jamʿiyya ʿalā maḥw al-ummiyya fī al-qurā.} — The association works for literacy in the villages.
  \item \emph{Recall:} say the Arabic for upbringing, then the Arabic for a kindergarten, then say \emph{maḥw al-ummiyya} again
\end{itemize}

\begin{tcolorbox}[breakable,colback=teal!4,colframe=teal!35!black,title={Before you move on}]
What does \ar{محو} \ar{الأمية} mean? (\textquotedblleft{}Literacy\textquotedblright{}.) Say it once more.
\end{tcolorbox}
\section[\texorpdfstring{\ar{حضور} (ḥuḍūr)}{حضور (ḥuḍūr)}]{\ar{حضور} (ḥuḍūr) — attendance}

\label{lesson:AR-C398-hudur}



\begin{quote}
Before the new one: say the Arabic for a kindergarten, then the Arabic for literacy.
\end{quote}

\subsection*{You'll want to know: \ar{حضور}}
\textbf{\ar{حضور}} — \emph{ḥuḍūr} — \textquotedblleft{}attendance\textquotedblright{}.

\begin{quote}
\ar{الحضور} \ar{إلزامي} \ar{في} \ar{هذا} \ar{الصف}.
\end{quote}

\begin{quote}
\emph{al-ḥuḍūr ilzāmī fī hādhā aṣ-ṣaff.}
\end{quote}

\begin{quote}
Attendance is compulsory in this class.
\end{quote}

\begin{cousinweb}
\textbf{\ar{حضور}} (\emph{ḥuḍūr}) is built on the native Arabic root \emph{ḥ-ḍ-r}, \textquotedblleft{}to be present\textquotedblright{} (Lane). \emph{Ḥuḍūr} is a verbal noun of \emph{ḥaḍara}, to be present: attendance.
\end{cousinweb}

\begin{grammarlens}[title={What you've built}]
One more word for giving reasons and opinions.
\end{grammarlens}

\subsection*{Your turn}
\begin{itemize}
\raggedright
  \item \emph{Say it:} \emph{ḥuḍūr}
  \item \emph{Say it:} \emph{ḥuḍūr}, once more
  \item \emph{Say it:} \emph{al-ḥuḍūr ilzāmī fī hādhā aṣ-ṣaff.} — Attendance is compulsory in this class.
  \item \emph{Recall:} say the Arabic for a kindergarten, then the Arabic for literacy, then say \emph{ḥuḍūr} again
\end{itemize}

\begin{tcolorbox}[breakable,colback=teal!4,colframe=teal!35!black,title={Before you move on}]
What does \ar{حضور} mean? (\textquotedblleft{}Attendance\textquotedblright{}.) Say it once more.
\end{tcolorbox}
